\documentclass[../../main.tex]{subfiles}
\begin{document}

{\bf [解答]}

\begin{figure}[htb]
  \centering
  \begin{tikzpicture}[x=2.5cm, y=3.0cm, >=stealth, every node/.style={font=\small}]
    % 座標軸
    \draw[->, thick, gray!80!black] (-0.7, 0) -- (2.6, 0) node[right] {$x$};
    \draw[->, thick, gray!80!black] (0, -0.6) -- (0, 1.2) node[above] {$y$};
    \node[below left=1pt, font=\footnotesize] at (0, 0) {$O$};

    % 原点での接線 y = x
    \draw[thick, dashed, teal!85!black, domain=-0.4:1.1] plot (\x, {\x})
    node[above left=1pt, font=\footnotesize, text=teal!85!black] {$y = x$};

    % 放物線 y = ax^2 + x (a = -0.5)
    \draw[line width=1.3pt, blue!85!black, domain=-0.4:2.3, samples=80]
    plot (\x, {-0.5*\x*\x + \x})
    node[below right=2pt, font=\footnotesize, text=blue!85!black] {$y = ax^2 + x$};

    % 代表点 (u, v): u = 1.4, v = -0.5*(1.4)^2 + 1.4 = 0.42
    \coordinate (P) at (1.4, 0.42);
    % 接線の傾き f'(u) = 2*(-0.5)*1.4 + 1 = -0.4
    % 接線: y - 0.42 = -0.4*(x - 1.4) => y = -0.4*x + 0.98
    \draw[line width=1.1pt, red!85!black, domain=0.4:2.3]
    plot (\x, {-0.4*\x + 0.98})
    node[right, font=\footnotesize, text=red!85!black] {接線};

    % 垂線 (u, v) -> x軸, y軸
    \draw[densely dotted, thick, gray!70!black] (P) -- (1.4, 0) node[below, font=\footnotesize] {$u$};
    \draw[densely dotted, thick, gray!70!black] (P) -- (0, 0.42) node[left, font=\footnotesize] {$v$};

    % 点 P(u, v)
    \node[circle, fill=red!85!black, inner sep=1.4pt] at (P) {};
    \node[above right=2pt, font=\footnotesize, text=red!85!black, fill=white, inner sep=0.8pt] at (P) {$(u, v)$};
  \end{tikzpicture}
  \caption{放物線 $y = ax^2 + x$ と点 $(u, v)$ における接線}
  \label{fig:parabola_tangent}
\end{figure}

題意の放物線を$y=f(x)=ax^2+bx$とおく$(\because f(0)=0)$．ただし$a\not=0$である．
$f'(0)=1$だから$b=1$である．故に
\begin{align}
  &f(x)=ax^2+x \\
  \therefore &f'(x)=2ax+1
\end{align}
となる．また$(u,v)$が放物線上にあることから$a=\dfrac{v-u}{u^2}$である．
故にグラフ上の点$(u,f(u))$での接線の傾きは
\begin{align}
  f'(u)&=2ua+1 \\
  &=2u\dfrac{v-u}{u^2}+1 \\
  &=\dfrac{2v-u}{u}\dots\text{(答)}
\end{align}
である．
\newpage
\end{document}
